\chapteropen{02}

\begin{cele}
\item nazwać swoją główną barierę czytania zamiast ogólnego \glqq nie mam czasu\grqq{},
\item odróżnić brak czasu od braku wyboru i sprawdzić tę różnicę w kalendarzu, a nie w samopoczuciu,
\item przeprowadzić autodiagnozę w pięciu pytaniach i odczytać jej wynik bez oceniania siebie,
\item wskazać rozdział, który rozwiązuje każdą z pięciu barier.
\end{cele}

\section{Pięć barier, nie jedna}

Zdanie \glqq nie mam czasu\grqq{} jest wygodne, bo nie zobowiązuje do niczego. Nie wskazuje, co zmienić, i nie daje się sprawdzić. Ludzie, którzy kupują książki i ich nie czytają, zwykle nie mają jednej przeszkody, tylko kilka, które maskują się tym samym zdaniem. W tym rozdziale wyróżniam pięć barier: brak czasu i energii, rozproszenie uwagi przez ekran, zły wybór książek, perfekcjonizm połączony z poczuciem winy oraz brak miejsca i rytuału. Każda wymaga innego działania. Dziesięć dodatkowych minut dziennie nie pomoże komuś, kto wybiera książki, które go nudzą, i nie pomoże komuś, kto kładzie się spać z telefonem w ręku.

Barier nie należy traktować jako cech charakteru. To opis mechanizmu, który działa w konkretnym tygodniu. Ten sam człowiek może mieć w marcu barierę energii po zmianie pracy, a w sierpniu barierę wyboru, bo skończył serię, która go trzymała.

\section{Czas, którego nie ma, i energia, której nie ma}

Najprostsza część diagnozy to arytmetyka, i dlatego warto ją zrobić przed resztą. Czytanie kosztuje minuty, a rok ma 365 dni. Piętnaście minut dziennie daje 5475 minut rocznie, czyli około 91 godzin. Typową książkę przyjmujemy jako \val{typical.words} słów. Przy 250 słowach na minutę zajmuje ona około \val{typical.hours} godziny. Stąd wynika, że przy 10 minutach dziennie czytasz w ciągu roku około \val{books.10.250} książek, a przy 30 minutach dziennie i 200 słowach na minutę --- \val{books.30.200} książki.

Z tej arytmetyki wynika coś, czego zdanie \glqq nie mam czasu\grqq{} nie mówi: liczbę, której potrzebujesz. Żeby przeczytać 12 książek w roku, potrzebujesz średnio \val{need.12} minuty dziennie. Na 24 książki wystarczy \val{need.24} minut, a na 52 --- \val{need.52} minuty. To nie są wielkie liczby. Jeśli nie możesz znaleźć nawet dziesięciu minut, problemem nie jest czas w sensie zegara, tylko to, dokąd on wypływa.

Druga bariera nie ma nic wspólnego z zegarem. Jest nią energia. Po pracy, przy dzieciach, po dyżurze często zostaje niewiele: można mieć godzinę wolnego i nie mieć na nią siły. Czytanie wymaga utrzymania uwagi na zdaniach, których tempa nie kontrolujesz, a to kosztuje więcej, niż się zdaje. Ludzie w takim stanie często obwiniają się o lenistwo, tymczasem dają książce to, co zostało po dniu. Gdy energia jest niska, dobrą reakcją nie jest wymuszanie godziny, tylko obniżenie wymagań: krótsza sesja, łatwiejsza książka, czytanie w kolejce. Sen wpływa na zdolność skupienia. Ogólne zalecenia dotyczące snu dla dorosłych mieszczą się w zakresie, który \verify{zwykle podaje się jako 7--9 godzin na dobę}; to, ile potrzebujesz konkretnie ty, sprawdzisz tylko na sobie.

\begin{przyklad}[Arytmetyka zamiast wyrzutów]
Wyobraź sobie wymyśloną Martę, która dojeżdża codziennie pociągiem. Dwie drogi po 25 minut dają 50 minut dziennie, czyli 250 minut tygodniowo przy pięciu dniach pracy. Gdyby czytała całą podróż przez 52 tygodnie, miałaby 13\,000 minut w roku. Przy 250 słowach na minutę to 3\,250\,000 słów, czyli około 40 książek typowej długości. Marta zwykle nie czyta w pociągu, bo ma telefon w dłoni. Jej bariera nie jest czasem. Jest nią sposób, w jaki spędza czas dojazdu. To przenosi problem z kalendarza na nawyk, a nawyki zmienia się inaczej niż grafik.
\end{przyklad}

\section{Ekran, który zjada uwagę}

Trzecia bariera to rozproszenie, i tu łatwo o nieprecyzyjność. Ekran nie zabiera nam czasu tak jak praca. Zabiera uwagę, która była do dyspozycji podczas czytania. Różnica jest subtelna: możesz mieć dwadzieścia minut wolnego, spędzić je z książką otwartą na tej samej stronie i nie pamiętać początku akapitu, bo co chwilę wracała myśl o powiadomieniu.

Badania nad uwagą opisują koszt przełączania zadań i wpływ urządzenia leżącego w zasięgu wzroku. Wnioski są mocniejsze, niż chciałby sam czytelnik: \verify{sama obecność telefonu w polu widzenia obniża wynik w testach uwagi, nawet gdy nie jest używany}. Traktuj to jako hipotezę do sprawdzenia na sobie, nie jako wyrok.

Praktyczny wniosek nie brzmi \glqq wyrzuć telefon\grqq{}. Brzmi: zmień kolejność. Jeśli czytanie zaczyna się po przejrzeniu skrzynki, sesja czytania zaczyna się już w rozproszeniu. Jeśli zaczyna się od książki, a telefon zostaje w innym pokoju, pierwsze minuty wyglądają inaczej. Ta zamiana, a nie siła woli, jest tym, co zwykle działa. Warto też wiedzieć, że przerwy w czytaniu nie są same w sobie szkodliwe. Szkodliwe jest to, co wypełnia przerwę: każde sprawdzenie skrzynki kończy się dziesięcioma minutami na sprawach innych niż książka, a pierwsza strona po powrocie trzeba czytać od nowa. Dlatego dobrym miernikiem nie jest długość sesji, tylko to, ile pierwszych minut sesji należało do tekstu. Rozdział 3 pokazuje, jak zamienić ekran na książkę w konkretnych oknach dnia; rozdział 7 mówi, co robić, gdy uwaga i tak się rozjeżdża, a tekst jest trudny.

\begin{falle}[Czytanie w telefonie jako \glqq czytanie w drodze\grqq{}]
Czytnik lub aplikacja w telefonie wyglądają jak rozwiązanie, bo książka jest zawsze pod ręką. W praktyce ta sama aplikacja pokazuje powiadomienia, zakładki i przycisk prowadzący do sklepu. Zanim przeczytasz rozdział, już wiesz, że możesz go przerwać, a to zmienia sposób czytania. Jeśli chcesz używać telefonu, ustaw tryb skupienia blokujący powiadomienia i sprawdź przez dwa tygodnie, czy liczba przeczytanych stron się zmieniła. Liczba stron to dana, której możesz użyć; \glqq wydaje mi się\grqq{} nie jest daną.
\end{falle}

\section{Zły wybór, perfekcjonizm i poczucie winy}

Czwarta bariera bywa mylona z lenistwem. Książka leży na stoliku od trzech miesięcy. Każdego wieczoru wzrok na niej trafia i odwraca się. Łatwo wtedy uznać, że to ty jesteś problemem, a może problemem jest ta konkretna książka, którą wybrałeś, bo wypadało.

Zły wybór ma kilka postaci: książka za trudna na obecny poziom uwagi, książka za długa na energię, którą masz, albo książka, która była czyjąś rekomendacją, a nie twoją potrzebą. Lekarstwo jest proste w opisie i trudne w praktyce: odłożyć książkę bez wyrzutów, według ustalonej wcześniej reguły. Rozdział 4 proponuje zasadę pięćdziesięciu stron: czytasz do pięćdziesiątej strony, a potem decydujesz, czy książka jest warta dalszego czasu. Liczba pięćdziesiąt nie wynika z badań. To reguła robocza, którą łatwo sprawdzić i którą można zmienić.

Perfekcjonizm działa w przeciwnym kierunku. Zakłada, że skoro czytasz, to powinieneś czytać dużo, uważnie, z notatkami i w odpowiednim tempie. Każde odstępstwo od tego obrazu zamienia się w poczucie winy, a poczucie winy zniechęca do kolejnej sesji. Ktoś, kto zaczyna rok od planu na 60 książek i w marcu ma cztery, często przestaje czytać w ogóle, bo zaległość wydaje się zbyt duża, żeby ją nadrobić. Korekta jest prosta: plan liczy się w minutach dziennych, a nie w tytułach, a zaległość zamyka się rewizją, nie karą. Rozdział 10 pokazuje, jak zrobić taką rewizję w grudniu.

\begin{zasada}[Odłożenie też jest czytaniem]
W tej książce przyjmujemy regułę roboczą: książka odłożona po rozsądnej próbie nie jest porażką czytelniczą. To decyzja o ograniczonym czasie, a nie dowód słabej woli. Wystarczy zapisać datę i powód, żeby odróżnić decyzję od zapomnienia.
\end{zasada}

\section{Miejsce i rytuał}

Piąta bariera wygląda najmniej poważnie, a bywa najtrwalsza. Nie masz miejsca, w którym czytanie ma sens: kanapa jest zajęta, łóżko służy do snu, biurko do pracy. Każdego wieczoru musisz podjąć decyzję, gdzie usiąść, i za każdym razem ta decyzja kosztuje. Rytuał zmniejsza ten koszt. Książka leży w tym samym miejscu, lampka zapala się o tej samej porze, a pierwsza czynność po zamknięciu laptopa to otwarcie książki. Brzmi to banalnie, ale działa, bo zamienia decyzję w odruch. Rozdział 6 mówi o wyzwalaczach, miejscu i planie awaryjnym na dni, w których nic nie wychodzi; rozdział 3 o oknach, które mają ustalony początek.

Osobnego pokoju nie trzeba. Wystarczy, że miejsce jest stałe i że da się w nim zacząć w ciągu minuty.

\section{Autodiagnoza w pięciu pytaniach}

Odpowiedz na pięć pytań, najlepiej na kartce. Odpowiedzi nie oceniaj; zapisz je.

\begin{enumerate}
\item Ile minut dziennie czytałeś w ostatnim tygodniu? Liczba ma pochodzić z zegarka lub kalendarza, nie z pamięci.
\item Ile książek zacząłeś w ostatnim roku i ile z nich skończyłeś?
\item Kiedy ostatnio sięgnąłeś po książkę zamiast po telefon i ile minut to trwało?
\item Co robisz w pierwszych piętnastu minutach wolnego czasu, zanim zdecydujesz, czym się zająć?
\item Jak się czujesz, gdy nie czytasz: winny, obojętny, zmęczony, ulżony?
\end{enumerate}

Interpretacja. Jeśli liczba z pierwszego pytania jest zgodna z tym, czego chcesz, a bariera leży gdzie indziej, czytanie nie jest problemem czasu. Jeśli liczba jest wyraźnie niższa niż zakładałeś, a w czwartym pytaniu pojawia się telefon, bariera jest w nawyku i w uwadze (rozdziały 3, 6 i 7). Jeśli wiele książek zostało zaczętych i nie skończonych bez decyzji, bariera leży w wyborze (rozdział 4). Jeśli w piątym pytaniu dominuje wina, bariera jest emocjonalna, i wtedy plan z mniejszą liczbą obietnic ma większą szansę niż plan z większą. Te odczytania są moją oceną, nie progami z badań. Możesz je przyjąć albo odrzucić, patrząc na własne liczby.

\begin{uwaga}[Nie mam czasu czy nie wybieram?]
Tę różnicę da się sprawdzić. Jeśli w kalendarzu jest czas przeznaczony na rzeczy mniej ważne dla ciebie, to nie brakuje czasu, tylko go nie wybierasz. Jeśli kalendarz jest naprawdę zapchany, pytanie brzmi inaczej: co przestajesz robić, żeby czytać. Obie odpowiedzi są uczciwe, ale wymagają różnych działań.
\end{uwaga}

\section{Koszty czytania i mapa barier}

Ta książka zakłada, że czytanie jest warte czasu. Nie każdy musi tak uważać. Czytanie kosztuje minuty, które mógłbyś przeznaczyć na naukę, rodzinę, sen, sport albo na nic. Jeśli wybierasz te inne rzeczy świadomie, nie ma w tym porażki. Książka ma pomóc ci przestać czytać z poczuciem winy albo ze złudzeniem, że to ty nie pasujesz do czytania. Nie ma wymuszać liczby na okładce.

Koszty są też konkretne. Książki kosztują pieniądze, a biblioteka wymaga dojazdu. Audiobook jest oszczędny w czasie, ale wymaga abonamentu i innej uwagi, co omawia rozdział 5. Czytanie zajmuje czas, który nie służy niczemu innemu, i nie zawsze jest to czas, którego wcześniej odmawiałeś sobie z nawyku.

Poniższa tabela zamyka diagnozę. Każda bariera ma rozdział, który ją rozwiązuje, i pytanie, które pozwala ją rozpoznać.

\begin{center}
\begin{tabular}{@{}p{0.27\linewidth}p{0.40\linewidth}p{0.22\linewidth}@{}}
\toprule
Bariera & Pytanie diagnostyczne & Rozdział \\
\midrule
Czas i energia & Ile minut dziennie naprawdę czytasz? & 1, 3, 6 \\
Ekran i uwaga & Co robisz w pierwszych minutach sesji? & 3, 5, 7 \\
Zły wybór & Ile książek odłożyłeś bez decyzji? & 4 \\
Perfekcjonizm i wina & Co czujesz, gdy nie czytasz? & 4, 10 \\
Miejsce i rytuał & Gdzie i o której godzinie zaczynasz? & 6 \\
\bottomrule
\end{tabular}
\end{center}

\begin{cwiczenia}
\exitem{Przez siedem dni zapisuj, ile minut dziennie naprawdę czytasz. Podaj średnią.}{Średnia z siedmiu zapisanych wartości w minutach. Dzień bez czytania też wpisz jako zero: zero jest daną, a nie brakiem danej.}
\exitem{Oblicz, ile minut dziennie potrzebujesz na 24 książki w roku, przy typowej książce 80\,000 słów i 250 słowach na minutę.}{$24 \times 80\,000 / 250 = 7\,680$ minut rocznie; $7\,680 / 365 \approx 21$ minut dziennie, czyli \val{need.24} minut.}
\exitem{Czytasz 200 słów na minutę. Ile książek typowej długości zmieści się w roku przy 30 minutach dziennie?}{$30 \times 365 \times 200 = 2\,190\,000$ słów; po podzieleniu przez 80\,000 wychodzi około 27,4, czyli \val{books.30.200} książki.}
\exitem{Przypomnij sobie ostatnią sesję, w której nie czytałeś mimo wolnego czasu. Co zrobiłeś jako pierwsze?}{Odpowiedź konkretna: nazwa aplikacji albo czynność. Jeśli była to czynność inna niż czytanie, bariera leży w zamianie, rozdział 3.}
\exitem{Podaj jedną książkę, którą odłożyłeś w ostatnim roku. Czy odłożenie było decyzją?}{Tak lub nie, z uzasadnieniem. Jeśli nie była to decyzja, tylko zapomnienie, zabrakło reguły, rozdział 4.}
\exitem{Napisz zdanie zaczynające się od \glqq nie czytam, bo\grqq{}, a potem przepisz je tak, by zaczynało się od \glqq nie wybieram czytania, kiedy\grqq{}.}{Przykład: \glqq nie czytam, bo nie mam czasu\grqq{} staje się \glqq nie wybieram czytania wieczorem, kiedy oglądam serial\grqq{}. Zamiana nie jest pocieszeniem; zmienia miejsce, w którym szukasz przyczyny.}
\end{cwiczenia}

\zadanie{Przez siedem dni zapisuj każdego dnia dwie rzeczy: minuty czytania i to, czym zastąpiłeś pierwsze piętnaście minut wolnego czasu. Miara sukcesu: po tygodniu masz siedem liczb, średnią i jedno zdanie, w którym nazywasz główną barierę po imieniu, z numerem rozdziału z tabeli.}
